% Clase 19 --- Catalogo intuitivo de comportamientos en un punto (§3).
% El docente hace estos cuatro dibujos antes de los ejemplos formales: "en
% general, cuando ven la gráfica, ya pueden predecir si tiene límite o si no
% tiene límite". Incluye sen(1/x) ("una función que oscila demasiado rápido")
% y la función signo ("ambos límites laterales son distintos").
\begin{center}
\begin{tikzpicture}
\begin{axis}[calcfig, width=6.4cm, height=4.2cm, axis lines=middle,
    xmin=-1.6, xmax=1.6, ymin=-0.6, ymax=2.15, xtick=\empty, ytick=\empty,
    xlabel={}, ylabel={}, title={\small (a) continua},
    title style={font=\small, yshift=-2pt}]
  \addplot[domain=-1.5:1.5, samples=60] {0.9 + 0.45*x};
  \node[figptocerrado] at (axis cs:0,0.9) {};
\end{axis}
\end{tikzpicture}\hspace{0.6em}
\begin{tikzpicture}
\begin{axis}[calcfig, width=6.4cm, height=4.2cm, axis lines=middle,
    xmin=-1.6, xmax=1.6, ymin=-0.6, ymax=2.15, xtick=\empty, ytick=\empty,
    xlabel={}, ylabel={}, title={\small (b) salto: $L \neq f(x_0)$},
    title style={font=\small, yshift=-2pt}]
  \addplot[domain=-1.5:-0.03, samples=40] {0.9 + 0.45*x};
  \addplot[domain=0.03:1.5, samples=40, forget plot] {0.9 + 0.45*x};
  \node[figptoabierto] at (axis cs:0,0.9) {};
  \node[figpto, fill=figred, minimum size=5pt] at (axis cs:0,1.62) {};
\end{axis}
\end{tikzpicture}\\[1em]
\begin{tikzpicture}
\begin{axis}[calcfig, width=6.4cm, height=4.2cm, axis lines=middle,
    xmin=-0.62, xmax=0.62, ymin=-1.5, ymax=1.5, xtick=\empty, ytick=\empty,
    xlabel={}, ylabel={}, title={\small (c) oscila: $\sin(1/x)$},
    title style={font=\small, yshift=-2pt}]
  \addplot[domain=0.028:0.6, samples=900] {sin(deg(1/x))};
  \addplot[domain=-0.6:-0.028, samples=900, forget plot] {sin(deg(1/x))};
\end{axis}
\end{tikzpicture}\hspace{0.6em}
\begin{tikzpicture}
\begin{axis}[calcfig, width=6.4cm, height=4.2cm, axis lines=middle,
    xmin=-1.6, xmax=1.6, ymin=-1.6, ymax=1.6, xtick=\empty, ytick=\empty,
    xlabel={}, ylabel={}, title={\small (d) laterales distintos},
    title style={font=\small, yshift=-2pt}]
  \addplot[domain=-1.5:-0.02, samples=2] {-1};
  \addplot[domain=0.02:1.5, samples=2, forget plot] {1};
  \node[figptoabierto] at (axis cs:0,-1) {};
  \node[figptoabierto] at (axis cs:0,1) {};
  \node[figptocerrado] at (axis cs:0,0) {};
\end{axis}
\end{tikzpicture}\\[0.5em]
{\small\itshape Los cuatro casos del pizarrón. En (a) hay límite y coincide
con el valor: es continua. En (b) hay límite pero $f(x_0)$ está en otro lado.
En (c) no hay límite en $0$ porque la función oscila cada vez más rápido (el
valor que se le asigne en $0$ es irrelevante). En (d) los dos límites laterales
existen pero difieren, así que tampoco hay límite.}
\end{center}
